\documentclass[10pt,a4paper]{amsart}
\usepackage[margin=19mm]{geometry}
\usepackage{amsmath,amssymb,mathtools}
\usepackage[scr=rsfs,cal=euler]{mathalfa}
\usepackage{xfrac}
\usepackage[unicode,hidelinks,bookmarks=false]{hyperref}
\newcommand{\ie}{i.e.\ }
\newcommand{\cf}{cf.\ }
\newcommand{\resp}{respectively\ }
\newcommand{\Subsection}{\subsection}
\providecommand{\nobreakdash}{\nobreak\hbox{-}}
\setlength{\emergencystretch}{2em}
\multlinegap0pt
\usepackage{bm}
\usepackage{bbm}
\usepackage{dsfont}
\usepackage{xspace}
\usepackage{cancel}
\usepackage{mathtools}
\usepackage[shortlabels]{enumitem}
\usepackage{amsthm}
\makeatletter
\@ifundefined{Theorem}{\newtheorem{Theorem}{Theorem}}{}
\@ifundefined{Lemma}{\newtheorem{Lemma}[Theorem]{Lemma}}{}
\@ifundefined{Proposition}{\newtheorem{Proposition}[Theorem]{Proposition}}{}
\@ifundefined{Remark}{\newtheorem{Remark}[Theorem]{Remark}}{}
\makeatother
\newcommand{\J}{\mathfrak{J}}
\newcommand{\N}{\mathcal{N}}
\newcommand{\U}{\mathcal{U}}
\newcommand{\g}{\mathfrak{g}}
\newcommand{\image}{\operatorname{Im}}
\newcommand{\m}{\mathfrak{m}}
\newcommand{\orb}{\mathfrak{O}}
\newcommand{\reg}{{\operatorname{reg}}}
\title[\'Etale geometry of closures of Jordan classes]{\'Etale geometry of closures of Jordan classes}
\author{Filippo Ambrosio}
\date{Source: arXiv:2411.07748v2 (CC BY 4.0)}
\begin{document}
\maketitle

\noindent\emph{Source and licence.} \'Etale geometry of closures of Jordan classes. Version arXiv:2411.07748v2 (CC BY 4.0), distributed under the Creative Commons Attribution 4.0 International licence, as recorded in the arXiv licence metadata for this exact version and shown on the abstract page https://arxiv.org/abs/2411.07748v2. Full rights notice: licenses/arxiv-2411.07748-CC-BY-4.0.txt.\par\smallskip

\noindent\textit{Editorial scope.} Counted unit: the Proposition whose source label is \texttt{prop\_respect} in the article (source0.tex lines 718--725, source label \texttt{prop\_respect}), delivered together with the article's own complete original proof of it (source0.tex lines 726--762, one \texttt{proof} environment of 37 lines). No mathematics is written, repaired or re-derived: statement and proof are the article's own text line for line. Required context retained uncounted: incl\_1 (lines 672--675); lem\_ss (lines 677--680); rk\_etale (lines 711--715). Every definition, display and label of those lines is retained unchanged. Omitted from this package and not redistributed: 1--671, 676--676, 681--710, 716--717, 763--910. Nothing inside a retained excerpt is omitted or altered: every retained body is delivered exactly as the article's lines are written, so no omission receipt of any kind accompanies this package. Citations: the retained excerpts carry no citation command at all, so no bibliography entry of the article is transcribed and no reference is redistributed. Reference adaptation: none was needed and none was made; every reference inside the delivered unit resolves inside it, so no unresolved reference or citation is produced. Printed slips kept literal: no printed word, symbol, hypothesis or number of the counted statement or of its proof was altered. Layout and class: the article's own class is \texttt{amsart} with packages including fullpage, amscd, bbm, amsaddr, amssymb, mathabx, mathrsfs, wasysym, tikz, tikz-cd, extpfeil, amsmath, enumitem, geometry; the wrapper is the lane's pinned \texttt{amsart} editorial wrapper at 10pt on A4, which restates the theorem-like environments the retained text needs and adds the article's own macro definitions that the retained text needs (\texttt{\textbackslash J}, \texttt{\textbackslash N}, \texttt{\textbackslash U}, \texttt{\textbackslash g}, \texttt{\textbackslash image}, \texttt{\textbackslash m}, \texttt{\textbackslash orb}, \texttt{\textbackslash reg}), with extra wrapper packages (bm, bbm, dsfont, xspace, cancel, mathtools, enumitem). The delivered numbering is the unit's own. Assets: no retained span contains a figure environment, a graphics-inclusion command, a PDF-page inclusion, a table or a TikZ picture; no image file is copied into this package, the package declares no asset dependency and no image file is redistributed with it. Licence: version arXiv:2411.07748v2 of this article is distributed under the Creative Commons Attribution 4.0 International licence, as recorded in the arXiv licence metadata for this exact version and shown on the abstract page https://arxiv.org/abs/2411.07748v2. A full rights notice is delivered at licenses/arxiv-2411.07748-CC-BY-4.0.txt. Characters: every retained excerpt is pure ASCII, so the delivered engine, XeLaTeX with its default Latin Modern text font, reports no missing character.
 \begin{align}
 G_h \subset G_{\lambda(h)}, \label{incl_1}\\
 \lambda(G_h) \subset \g_h. \label{incl_2}
 \end{align}

\begin{Lemma} \label{lem_ss}
A $G$-equivariant morphism $\lambda \colon G \to \g$ maps semisimple elements of $G$ to semisimple elements of $\g$.
Moreover, if $su$ is the Jordan decomposition of $g \in G$, then $\lambda(s) + (\lambda(g) - \lambda(s))$ is the Jordan decomposition of  $\lambda (g) \in \g$.
\end{Lemma}

\begin{Remark}\label{rk_etale}
For any element $h \in  U_\lambda$, thanks to  \eqref{incl_1} and since  \'etale maps preserve dimensions, we have $G_h^\circ = G_{\lambda(h)}^\circ$.
If moreover the subgroup $G_{\lambda(h)}$ is itself connected, then $G_h^\circ =G_h =   G_{\lambda(h)}$ (in particular, if $h$ is semisimple and $p$ is not torsion for $G$.)
Under separability assumptions (for example if $h$ is semisimple), we also have $\g_h = \g_{\lambda(h)}$.
\end{Remark}

\begin{Proposition} \label{prop_respect}
Let $\lambda \colon G \to \g$ be a log-like map and let $su \in U_\lambda$.
The following assertions hold.
\begin{enumerate}[label=(\arabic*)]
\item $\lambda( J(su) \cap U_\lambda) \subset \J(\lambda(su)) \cap V_\lambda$.
\item If the identity element $e \in \overline{J(su)}$, then $\J(\lambda(su)) = \J(\lambda(s) + \lambda(u))$.
\end{enumerate}
\end{Proposition}

\begin{proof}
We prove (1). For $su \in G$ put $G_s^\circ =: M, Z(M) =: Z$ so that $\m = \g_s$.
It suffices to show that $\lambda(((Z^\circ s)^\reg \cap U_\lambda) u) \subset \J(\lambda(su))$.
For $z \in (Z^\circ s)^\reg \cap U_\lambda$, we  show that $\lambda(zu)$ is Jordan equivalent to $\lambda(su)$.
Recall that we know the Jordan decomposition of such elements thanks to Lemma  \ref{lem_ss}.

We first deal with the semisimple part.
We have $M = G_z^\circ$ and thanks to  Remark \ref{rk_etale} we derive also $M = G_{\lambda(s)}^\circ = G_{\lambda(z)}^\circ$.
In particular $\m = \g_{\lambda(z)} = \g_{\lambda(s)}$.

For the nilpotent part, we define  the $M$-equivariant morphism 
\begin{align*}
\nu \colon M \times Z^\circ s &\to \N(\m) \\
(g, z) & \mapsto g \cdot (\lambda(zu) - \lambda(z)) = \lambda(z gug^{-1}) - \lambda(z),
\end{align*}
which is well-defined by Lemma \ref{lem_ss}, in particular the Jordan decomposition $\lambda(zu) = \lambda(z) + \nu(e,z)$ holds in $\m$ for all $z \in Z^\circ s$.
Here $M$ acts on the domain of $\nu$ by left multiplication on the first factor and on the codomain by the adjoint action.
The closure $\overline{\image \nu}$ is an irreducible $M$-stable subset of $\N(\m)$, hence $\overline{\image \nu} = \overline{\orb}$ for some nilpotent $M$-orbit $\orb \subset \m$.
Fix $d := \dim \orb$, then $\orb = \overline{\image \nu} \cap \m_{(d)}$ and $M$-equivariance implies $\orb \subset \image \nu$ so that $\orb = \image \nu \cap \m_{(d)}$.
We prove that $\varnothing \neq \nu (M \times ( Z^\circ s)^\reg \cap U_\lambda )) \subset \m_{(d)}$.
We  have $s \in (Z^\circ s)^\reg \cap U_\lambda$, proving non-emptiness.
Let $z \in (Z^\circ s)^\reg \cap U_\lambda$.
Then the \'etale property of $\lambda$ on $U_\lambda$ implies: $$\dim G \cdot (\lambda(z) + \nu(e,z)) = \dim G \cdot (zu) = \dim G \cdot (su) = \dim G \cdot (\lambda(s) + \nu(e,s)).$$
This equality yields 
$\dim G_{\lambda(z)} \cdot \nu(e,z) = \dim G_{\lambda(s)} \cdot \nu(e,s)$ which in turn implies
$\dim M \cdot \nu(e,z) = \dim M \cdot \nu (e,s)$, whence
$M \cdot \nu(e,z) \subset \m_{(d')}$
with $d' = \dim G \cdot (su) - \dim G \cdot s = \dim M \cdot u$.
The set $(Z^\circ s)^\reg \cap U_\lambda$ is open in $Z^\circ s$ so $\nu(M \times ((Z^\circ s)^\reg \cap U_\lambda))$ is dense in $\image \nu$ and must meet $\orb$ nontrivially, which implies $\orb = \nu(M \times ((Z^\circ s)^\reg \cap U_\lambda))$ so that $d' = d$ and we conclude.
In particular $\J(\lambda(su)) = \J(\m, \orb)$.

To prove (2), recall that $J(su)$ closes at the identity if and only if $Z^\circ s = Z^\circ$.
We already observed in (1) that $d' = \dim M \cdot u$.
Hence $d' = \dim M \cdot \lambda(u)$ since $\lambda$ induces an isomorphism $\U(M) \to \N(\m)$.
Finally we have $\lambda(u) = \nu(e,e) \in \overline{\image \nu}$, yielding
$M \cdot \lambda(u) \subset \overline{\image \nu} \cap \m_{(d')} = \orb$, hence $\J(\lambda(su)) = \J(\lambda(s) + \lambda(u))$.
\end{proof}

\end{document}
